% ID: SEL-D0029-OA-E02
% ESTADO: validada. S030, cotejada con examen y criterios.
\textbf{a)} El dominio es $\mathbb R$: el denominador del primer
tramo no se anula para $x<3$. En $3$, el límite izquierdo es
$(-2)/(-1)=2$, igual a $-9+21-10=2$, valor y límite derecho.
Los dos tramos son continuos en sus interiores.
Sus derivadas son $1/(x-4)^2$ y $7-2x$. Las pendientes laterales
coinciden en $3$, ambas valen 1. Por tanto $f$ es continua
y derivable en todo $\mathbb R$.

\textbf{b)} El eje vertical se corta en $(0,5/4)$.
El primer tramo no tiene ceros en $x<3$.
En el segundo, $-x^2+7x-10=-(x-2)(x-5)$, y sólo $5$
pertenece al tramo. El corte horizontal es $(5,0)$.

\textbf{c)} A la izquierda, $f(x)=1-1/(x-4)$ tiende a $1$
cuando $x\to-\infty$: $y=1$ es asíntota horizontal por ese lado.
No hay asíntotas verticales porque $f$ es continua en todo
$\mathbb R$. A la derecha el término cuadrático domina, así que
$f(x)$ se aleja indefinidamente de cualquier recta.
